\documentclass[11pt]{article}

\usepackage{amsmath,amssymb}
\usepackage[margin=1.0in]{geometry}
\usepackage{bm}
\usepackage{booktabs}
\usepackage{array}
\usepackage{graphicx}
\usepackage{float}
\usepackage{enumitem}
\newcolumntype{P}[1]{>{\raggedright\arraybackslash}p{#1}}
\usepackage[colorlinks=true,linkcolor=blue,citecolor=blue,urlcolor=blue]{hyperref}
\urlstyle{tt}
\makeatletter\g@addto@macro\UrlBreaks{\do\_\do\-}\makeatother

\input{macros.inc}

\numberwithin{equation}{section}
\sloppy
\setlength{\emergencystretch}{3em}
\graphicspath{{figures/benchmarks/}}

% --- notation (follows quasi_static_models.tex and benchmarks.tex) ----------
\newcommand{\R}{\mathbb{R}}
\newcommand{\eps}{\varepsilon}
\newcommand{\dS}{\,\mathrm{d}S}
\newcommand{\dV}{\,\mathrm{d}V}
\newcommand{\vu}{\mathbf{u}}
\newcommand{\vw}{\mathbf{w}}
\newcommand{\vx}{\mathbf{x}}
\newcommand{\vt}{\mathbf{t}}
\newcommand{\vn}{\hat{\mathbf{n}}}
\newcommand{\rhat}{\hat{\mathbf{r}}}
\newcommand{\code}[1]{\texttt{#1}}
\newcommand{\Ylm}{Y_{lm}}
\newcommand{\gradone}{\nabla_{1}}
\newcommand{\bsig}{\bm{\sigma}}

% A tier's status line: \status{date}{text}. Update the date with the text.
\newcommand{\status}[2]{%
  \par\smallskip\noindent
  \fbox{\parbox{\dimexpr\linewidth-2\fboxsep-2\fboxrule\relax}{\small
  \textbf{Status, #1.} #2}}\par\medskip}

\title{The benchmark programme of \texttt{AdGIA}:\\
       a hierarchy of forward tests with \texttt{sl3dpp} and \texttt{pyslfp}}
\author{}
\date{Begun 21 October 2026}

\begin{document}
\maketitle

\begin{abstract}
\noindent
This note is the write-up of the benchmarks of \code{AdGIA}, organised
around the cross-code programme begun in October 2026 with the two
spectral codes \code{sl3dpp} and \code{pyslfp}. The programme climbs a
hierarchy: elastic spherical models first (non-gravitating, then
gravitating; solid, then fluid--solid), then viscoelastic loading without
the sea-level equation, and only then sea level, so that a disagreement is
found on the lowest rung able to show it. Part~\ref{part:programme} sets out
the programme, its tiers and their status, and the conventions every tier
shares: one model definition, one exchange format, tiered tolerances, and
the rules at degrees zero and one. Parts~\ref{part:elastic}--\ref{part:sea}
follow the hierarchy. Each tier has a section with a dated status line and
a fixed shape (the problem, the compared quantities, the references and the
legs, the results, the open points), so that tiers are added as they land
without reworking what is already written. The internal benchmark families
of \path{benchmarks/} --- benchmarks of a different kind, against exact or
independent references --- are placed where they belong in the hierarchy,
with their conclusions summarised; their detailed measurements stay in
\code{doc/benchmarks.tex}, which this note cites and does not repeat. At the
time of writing the first tier, E1 (non-gravitating elastic, solid
spherical models), has its exact reference and both spectral legs in place
and gated, and its finite-element leg, \code{AdGIA}'s own, has its first
measurements; its refinement study across the models is next.
\end{abstract}

\tableofcontents

%=============================================================================
\part{The programme}
\label{part:programme}
%=============================================================================

\section{Why a hierarchy}
\label{sec:why}

A forward calculation of glacial isostatic adjustment stacks several
layers: the elastic response of a layered body; gravity and its
perturbation; fluid regions; the rheology and its time stepping; lateral
structure; and the sea-level equation, with its shorelines and its
rotational feedback. Two codes that disagree on the complete problem may
disagree in any of these, and a comparison made only at the top says how
large the spread is, not where it comes from. The published GIA benchmarks
show both the value and the limit of comparisons made high in the stack.
Spada et al.\ \cite{Spada2011} compared eight codes on one incompressible
Maxwell model, largely in the Laplace domain, and deliberately left out
compressibility and the sea-level equation. Martinec et al.\
\cite{Martinec2018} compared ten codes on the sea-level equation, on one
incompressible, non-rotating radial model; the codes agreed to about
1.5\,\%, and spatial discretisation and Gibbs behaviour at the load margins
were named as the main sources of the differences.

This programme goes the other way. It breaks the problem down as far as it
goes and compares the codes at every rung, so that each new ingredient is
tested on top of rungs already known to agree. Its principles:
\begin{itemize}[leftmargin=*]
\item \emph{Hierarchy, not headline.} Elastic before viscoelastic before
  sea level; non-gravitating before gravitating; simple spherical models
  before complex ones; loading without the sea-level equation before
  loading with it.
\item \emph{A practical target.} A GIA code must give stable, reliable
  results relative to \emph{observational} errors. Agreement to machine
  precision is not the goal, although understanding why two codes differ
  is valuable, and where two codes solve the identical discrete problem
  their agreement is expected at that level and is used as a gate.
\item \emph{Exact references where they exist.} The non-gravitating
  piecewise-homogeneous models have closed-form solutions, which is why the
  hierarchy starts there. Semi-analytical references (modal sums, inverse
  Laplace routes such as AAA with Talbot, or Gaver--Stehfest) are useful but
  secondary; some are fragile, particularly those that need an inverse
  Laplace transform.
\item \emph{Two modes for \code{AdGIA}.} Low-resolution runs that fit on a
  laptop are the default gate; server campaigns push the resolution later.
  The benchmarks are not part of continuous integration, which holds each
  code's own tests.
\end{itemize}

\section{The three codes}
\label{sec:codes}

\begin{center}
\small
\begin{tabular}{@{}lP{6.2cm}P{6.3cm}@{}}
\toprule
code & method & role in the programme \\
\midrule
\code{AdGIA} & 3-D finite elements on MFEM: quasi-static, self-gravitating
(visco)elastic loading; fluid regions by several formulations; sea level
solved monolithically; aspherical bodies by relabelling &
the code under test at every tier; the only one for general 3-D elastic
structure \\
\code{sl3dpp} & C++; spectral elements (GLL) in radius, spherical
harmonics in angle; time-domain viscoelasticity by implicit Prony
stepping; laterally varying viscosity $\eta=\bar\eta(r)\,e^{\delta}$ &
a second implementation of \code{pyslfp}'s radial method, and the second
code with 3-D viscosity \\
\code{pyslfp} & Python, released; generalised Love numbers on a radial
GLL mesh; the sea-level equation pseudo-spectrally; a time-domain
viscoelastic module (new in v2.2.1) &
the elastic oracle, the most thoroughly tested code at the radial level;
the 1-D viscoelastic reference \\
\bottomrule
\end{tabular}
\end{center}

\noindent Three things are shared and shape the comparisons.
\code{sl3dpp}'s degree operator mirrors \code{pyslfp}'s discretisation, so
on the same model and the same radial mesh rule the two solve the identical
discrete problem, while \code{AdGIA} is an independent discretisation. All
three take their models, decks, radial meshes and units from
\code{planetmodel}. And all three implement the same exponential-trapezoid
step weights, so that a viscoelastic comparison isolates the spatial
discretisation; \code{pyslfp}'s generalised sea-level equation is the
template of the other two codes' treatments.

\section{Two kinds of benchmark}
\label{sec:kinds}

The cross-code tiers set the codes against one another and, where one
exists, against an exact solution, always on spherical models and always
through a small set of numbers. Beside them, \path{benchmarks/} holds
\code{AdGIA}'s \emph{internal families}: Love numbers against
\code{pyslfp}, the relabelling identities, the interface-shift derivatives,
the two-dimensional disc, the viscoelastic box, sphere, stepping and
Love-history sub-families, and the sea-level fingerprints. These test
against exact closed forms, against identities that hold by construction
(the discrete change of variables under a relabelling), against refinement,
or against \code{pyslfp} as an independent radial reference. They reach
what no cross-code tier can --- the aspherical machinery, the
formulation-against-formulation questions, the time integrators in
isolation --- and they stay. Parts~\ref{part:elastic}--\ref{part:sea} place
each family at the rung where it belongs and state what it has
established; \code{doc/benchmarks.tex} remains the archive of their
measurements, problems and figures up to October 2026.

\section{The tiers and their status}
\label{sec:tiers}

Table~\ref{tab:tiers} lists the tiers as they stand on 21 October 2026.
``Done'' means implemented,
measured and recorded in this note; ``in place'' means the ingredient
exists but has not yet been brought under the programme's conventions;
``pending'' means not started. The labels of the viscoelastic and
sea-level sub-tiers are provisional until those tiers are designed.
\begin{table}[tp]
\centering
\small
\begin{tabular}{@{}lP{3.6cm}P{2.2cm}P{2.8cm}P{4.2cm}@{}}
\toprule
tier & problem & codes & reference & status \\
\midrule
E1 (\S\ref{sec:e1}) & elastic, non-gravitating, solid spherical models &
all three & exact (piecewise-homogeneous propagator) &
reference done; \code{pyslfp} and \code{sl3dpp} legs done and gated;
\code{AdGIA}'s finite-element leg implemented, first measurements;
refinement study pending \\ \addlinespace
E2 (\S\ref{sec:e2}) & elastic, gravitating, solid models & all three &
\code{pyslfp} & begun: the spectral pair gated on the gravitating traction
numbers; \code{AdGIA}'s Love-number family in place, exchange bridges being
written; degree-one frame work pending \\ \addlinespace
E3 (\S\ref{sec:e3}) & elastic, gravitating, fluid--solid models & all
three & \code{pyslfp} & \code{AdGIA}'s Love-number family in place
(Dahlen at $l\ge1$, Maxwell core at every degree); cross-code pending \\ \addlinespace
V (\S\ref{sec:v}) & viscoelastic loading without the sea-level equation:
model ladder, gravity off/on, rotation off/on & all three & the spectral
codes; semi-analytical secondary & pending; \code{AdGIA}'s viscoelastic
family and the spectral codes' own gates in place \\ \addlinespace
V-lat (\S\ref{sec:v}) & laterally varying viscosity, radial elastic
structure, no sea level & \code{AdGIA}, \code{sl3dpp} & each other &
pending; needs public \code{sl3dpp} drivers \\ \addlinespace
S (\S\ref{sec:s}) & sea level: elastic fingerprints (rotation off/on,
shoreline migration), then viscoelastic GIA runs; the lateral pair & all
three; lateral pair \code{AdGIA}, \code{sl3dpp} & \code{pyslfp}; each other
& pending; \code{AdGIA}--\code{pyslfp} and \code{sl3dpp}--\code{pyslfp}
comparisons in place separately \\
\midrule
internal & 3-D elastic (relabelling, perturbation, disc); viscoelastic
(box, sphere, stepping, Love histories); sea level & \code{AdGIA} & exact,
identities, \code{pyslfp} & in place (\S\ref{sec:internal-elastic},
\S\ref{sec:internal-visco}, \S\ref{sec:internal-sea}) \\
\bottomrule
\end{tabular}
\caption{The tiers of the programme and their status on 21 October 2026.}
\label{tab:tiers}
\end{table}
Three-dimensional elastic structure has no cross-code tier, since only
\code{AdGIA} computes it; it is covered by the internal families
(convergence, formulation against formulation, the mapped identities). A
comparison with another group's 3-D finite-element code may come later.

\section{How this note grows}
\label{sec:grows}

Each tier is one section with a dated status line and the same shape:
\emph{the problem}, \emph{the compared quantities}, \emph{the references
and the legs} (what each code computes, where, and how it is gated),
\emph{results}, and \emph{open points}. A tier not yet started records what
has been decided for it and what already exists. Measurements are quoted
with their conditions --- model, degrees, mesh rule or resolution, and the
commits where they are recorded --- and a result superseded by a later run
is replaced, not appended to. Durable content from the planning notes
(\path{~/dev/planning/}, which are working documents and not permanent)
moves here as each tier lands; how to run each benchmark stays in the
READMEs of \path{benchmarks/}.

%-----------------------------------------------------------------------------
\section{Conventions}
\label{sec:conventions}
%-----------------------------------------------------------------------------

\subsection{Models: one definition}
\label{sec:models}

The models are the spherical ladder of \path{benchmarks/common/models.py}:
planetmodel \code{LayeredIsotropicElastic} models in SI, each layer
constant or polynomial in radius for $\rho$, $v_P$ and $v_S$, fluid layers
with $v_S=0$. The values are listed in \code{doc/benchmarks.tex} (``The
models''). The rungs used or planned by the tiers:
\begin{center}
\small
\begin{tabular}{@{}P{4.2cm}P{5.4cm}P{5.4cm}@{}}
\toprule
model & structure & tiers \\
\midrule
\code{homogeneous} & one solid layer & E1 (exact), E2, V \\
\code{homogeneous\_lithosphere} & the same body cut at 5700\,km, so that
the outer shell can stay elastic & E1 (exact: elastically the homogeneous
body, a welded interface with no contrast), V \\
\code{two\_solid} & two homogeneous solid layers, interface at 5000\,km &
E1 (exact), E2, V \\
\code{linear\_solid} & one solid layer, linear in radius & E2, V \\
\code{fluid\_core}, \code{aw\_core} & homogeneous fluid core (non-neutral),
and its neutrally stratified twin & E3, V \\
\code{inner\_core}, \code{stratified\_core}, \code{earth\_like},
\code{prem\_4}, \code{prem\_6} & inner cores, stratified cores, PREM-grade
structure & E3, V, S \\
\bottomrule
\end{tabular}
\end{center}

One model object feeds every code. \code{pyslfp} takes the planetmodel
objects directly. \code{sl3dpp} reads tabulated decks, written from the same
objects by \path{benchmarks/common/decks.py}: the elastic deck format
$r,\rho,v_{PV},v_{SV},v_{PH},v_{SH},\eta$ in SI (radius first; repeated
radii mark internal boundaries; $\eta$ the anisotropy parameter, one for
the isotropic ladder), eight knots per layer, every value written to full
precision (\code{\%.17g}). The deck is read back with a not-a-knot cubic
spline, which reproduces a polynomial of degree at most three exactly; every
ladder layer is such a polynomial except \code{aw\_core}'s closed-form
$v_P$, so for every other model the deck \emph{is} the model to rounding;
\code{decks.roundtrip\_error} measures the agreement, so that a generation
script can state it. \code{AdGIA} meshes the objects themselves
(\path{benchmarks/common/make_case.py}, through planetmodel and gmsh), as
its internal families do.

The radial mesh is part of the discrete problem on the spectral side.
\code{pyslfp}'s default is its graded mesh (\code{graded\_mesh}: five GLL
nodes per element, fine at the surface and widening with depth, so that
every degree to $l_{\max}$ is resolved at least as the uniform rule would
resolve it, with a few hundred elements). The uniform rule for a given
$l_{\max}$, elements no wider than $0.1a/(l_{\max}+1)$, is shared by the two
spectral codes
(\code{planetmodel}'s \code{RadialMesh} and \code{sl3dpp}'s
\code{RadialMesh::Create}); the same deck read on the same rule is the
identical discrete problem in both.

\subsection{Harmonics and the numbers of one degree}
\label{sec:harmonics}

On a spherically symmetric model the response to a forcing of one
spherical harmonic $\Ylm$ separates as
\begin{equation}
  \vu = U(r)\,\Ylm\,\rhat + V(r)\,\gradone\Ylm,
  \qquad \phi = \phi(r)\,\Ylm,
  \label{eq:separation}
\end{equation}
with $\gradone$ the surface gradient on the unit sphere. The surface values
of $U$, $V$ and $\phi$ per unit forcing depend on the degree $l$ alone: not
on the order $m$, and not on how the harmonics are normalised, since forcing
and response carry the same harmonic. The codes use real orthonormal
harmonics ($\int_{S^2}\Ylm^2\,\mathrm{d}\Omega=1$, Condon--Shortley phase).
The normalisation matters only where a code builds or analyses fields:
\code{AdGIA}'s finite-element legs load with a harmonic and analyse the
solution on the surface with the same harmonics
(\code{BoundaryHarmonicCoefficients}).

\subsection{Units}
\label{sec:units}

Everything exchanged is in SI. The codes solve in different internal units
--- \code{AdGIA} in the units of \code{models.scales} (length the outer
radius, density 5000\,kg\,m$^{-3}$, time chosen so that $G=1$),
\code{pyslfp} and \code{sl3dpp} in \code{planetmodel} \code{Scales} with
different time scales --- so each code converts its own output and the
comparison is anchored in SI. A shared non-dimensionalisation is an open
cross-code question (\S\ref{sec:open}).

\subsection{The exchange format}
\label{sec:exchange}

Each code reduces its answer to a small text file: provenance lines, named
scalars, and one table. \path{benchmarks/common/exchange.py} writes and
reads it, and the C++ drivers reproduce its layout by hand.
\begin{itemize}[leftmargin=*]
\item Leading \code{\#} lines are provenance, free text. By convention the
  first says what the file holds and one reads \code{code <name> commit
  <sha> date <iso-date>}.
\item Bare \code{name value...} lines are named scalars.
\item A \code{columns} line names the columns of the table that follows,
  one row per line, every value SI and written to round-trip precision
  (\code{\%.17g}, or shortest round-trip in C++).
\item \code{nan} marks an excluded entry (a degree the problem does not
  define).
\end{itemize}
The first lines of the analytic E1 reference for \code{homogeneous}, as
\path{reference.py} writes it (two long lines shortened here):
\begin{verbatim}
# E1 non-gravitating traction numbers: the exact piecewise-homogeneous ...
# code analytic commit 5f2a5a3 date 2026-10-21
# model homogeneous SI (m per Pa); NaN = excluded degree
radius 6371000
gravitating 0
columns l h_u l_u h_v l_v
0 -6.4713052310817677e-06 0 0 0
1 nan nan nan nan
2 -1.3348735500072317e-05 -1.9325902953678312e-06 -1.1595541772206983e-05 ...
\end{verbatim}
Files are named by tier, model and code (\path{e1_<model>_<code>.txt}). The
format is deliberately plain: the aim is to reduce each comparison to a few
numbers or a simple plot. A genuine 3-D field (the lateral-viscosity tier)
will be exchanged as netCDF instead.

\subsection{Comparison and the error measure}
\label{sec:measure}

A comparison reads two exchange files and never needs the codes that wrote
them. \path{benchmarks/elastic/nongrav/compare.py}, the first such script,
takes the first file as the reference, matches the degree column and the
shared columns by name, and reports for each channel the worst relative
difference over the degrees, measured against the \emph{largest reference
value of the degree} over the shared channels:
\begin{equation}
  e_c = \max_l
  \frac{\bigl|x_c^{\text{other}}(l)-x_c^{\text{ref}}(l)\bigr|}
       {\max_{c'}\bigl|x_{c'}^{\text{ref}}(l)\bigr|} .
  \label{eq:degreescale}
\end{equation}
A \code{nan} in either file excludes the entry, and \code{--tol} turns the
report into a pass/fail exit status. The degree scale keeps a channel that
passes through zero from dominating: on \code{homogeneous}, $l_u$ changes
sign between degrees 3 and 4 (\S\ref{sec:e1-closed}), and at degree 4 it is
$8\times10^{-4}$ of $h_u$, so a relative error on its own value would
measure nothing physical.

The comparison rules, recorded first in \code{sl3dpp}'s decision log, hold
throughout: compare in coefficient space or through grid-independent
functionals, never point-wise on matched grids; compare gauge-invariant
quantities only; anchor in SI.

\subsection{Tolerances in tiers}
\label{sec:tolerances}

What counts as agreement depends on what is being compared:
\begin{enumerate}[leftmargin=*,label=(\roman*)]
\item \emph{The identical discrete problem} --- \code{pyslfp} and
  \code{sl3dpp} on the same deck and mesh rule. Agreement is expected near
  machine precision and is a gate: $10^{-10}$ of the degree scale.
\item \emph{An exact reference against a spectral code} --- the code's own
  discretisation error, spectrally small and dependent on the mesh rule;
  about $10^{-11}$--$10^{-10}$ in E1 (\S\ref{sec:e1-results}).
\item \emph{A finite-element code against a reference} --- the
  discretisation level, established by refinement in $h$ and in order:
  judged by the levels and by convergence under refinement, not by a fixed
  threshold.
\end{enumerate}
The physical verdict is set by observational error: the codes must be
reliable against what can be observed, not against machine epsilon. The
yardsticks for particular observables are to be fixed when the sea-level
tier defines its observables (\S\ref{sec:open}).

\subsection{Degrees zero and one}
\label{sec:degrees01}

\paragraph{Degree zero.} \code{pyslfp} solves degree 0 with the full
compressible (Lagrangian) equations, fluid regions included as elastic
media with zero shear modulus. \code{sl3dpp}'s rewrite was believed to do
the same; for solid models this is now verified, since it reproduces
\code{pyslfp}'s degree-0 traction numbers on the identical discrete problem,
gravitating and not (\S\ref{sec:e1-results}); the fluid--solid case falls
to E3. In \code{AdGIA}, Dahlen's mixed formulation with a fluid region
discards the fluid's bulk modulus and leaves an arbitrary function at
degree 0, so it differs from the others by design there: it is not
benchmarked at $l=0$ except to show the difference. The Maxwell-core
treatment (\code{SetFluid}'s Maxwell relaxation) keeps the compressible
fluid and is the route benchmarked at $l=0$. On solid models, as in E1 and
E2, none of this arises.

\paragraph{Degree one.} Without gravity, a spheroidal surface traction of
degree 1 has a net force, which a free body cannot balance; the
non-gravitating tiers exclude the degree and write \code{nan}
(\S\ref{sec:e1-solvability}). With gravity, degree-1 quantities are
compared in the frame of the centre of mass of the body and its load, in
which the external potential has no degree-one part; \code{AdGIA}'s
Love-number driver shifts its solution into that frame
(\code{doc/benchmarks.tex}, ``Degree one: the centre-of-mass frame''), but
its agreement with \code{pyslfp}'s frame is still unconfirmed, and E2
settles it. The rotational (toroidal, rigid-rotation) part of degree 1 is
more a convention than a frame and is documented by each code separately;
where it matters, the comparison uses quantities invariant under rigid
motions, such as changes of length between points. Comparisons at
$l\ne1$ are not affected.

\subsection{Where things live}
\label{sec:where}

Everything cross-code --- models, decks, exchange, comparison scripts,
results and this note --- lives in \code{AdGIA}; \code{sl3dpp} and
\code{pyslfp} keep their own internal tests. The tiers get directories of
\path{benchmarks/} as they land: \path{elastic/nongrav/} (E1) exists and
\path{elastic/grav/} (E2, E3) is begun; the planned layout adds
\path{elastic/internal3d/}, \path{viscoelastic/spherical/},
\path{viscoelastic/lateral/}, \path{sea_level/elastic/} and
\path{sea_level/viscoelastic/}, beside the existing family directories,
which stay. The shared pieces are in \path{benchmarks/common/}
(\code{models.py}, \code{decks.py}, \code{exchange.py}, \code{analytic.py},
\code{outputs.py}). \code{sl3dpp}'s drivers live in the \code{sl3dpp}
repository, built with \code{-DSL3DPP\_BUILD\_BENCHMARKS=ON}, and are used
from here; scripts that call \code{pyslfp} run from the \code{pyslfp}
checkout, which owns its environment. Every script refuses to write into
the source tree (\path{common/outputs.py}).

\section{Open points across the programme}
\label{sec:open}

\begin{itemize}[leftmargin=*]
\item \emph{Observational yardsticks.} Which observables, and what error on
  each, sets the physical verdict at each tier.
\item \emph{The degree-one rotational convention.} To be stated in each
  code's own documentation and compared through invariants where needed.
\item \emph{A shared non-dimensionalisation}, if kernels are to be
  exchanged between codes; until then, SI.
\item \emph{Public \code{sl3dpp} drivers} for the load-driven
  (no sea-level) viscoelastic tiers and for the lateral-viscosity pair.
\item \emph{Automation.} The tiers are not yet stages of
  \path{benchmarks/campaign.py}; the laptop and server profiles will run
  them as they are added.
\end{itemize}

%=============================================================================
\part{Elastic spherical models}
\label{part:elastic}
%=============================================================================

\section{Tier E1: non-gravitating elastic, solid models}
\label{sec:e1}

\status{21 October 2026}{The exact reference (\path{common/analytic.py})
is done. The \code{pyslfp} and \code{sl3dpp} legs are done and gated: on the
identical discrete problem against each other, and against the exact
reference (\S\ref{sec:e1-results}). \code{AdGIA}'s finite-element leg,
\code{traction\_benchmark}, is implemented and has first measurements on
\code{two\_solid} at two resolutions; the refinement study across the
three models is pending (\S\ref{sec:e1-fe}).}

\subsection{The problem}
\label{sec:e1-problem}

A solid ball $B$ of radius $a$, spherically layered, with Lam\'e moduli
$\lambda(r),\mu(r)$, no gravity, no pre-stress and no body force, is held
in static equilibrium by a prescribed traction $\vt$ on its surface:
\begin{equation}
  \nabla\cdot\bsig = \mathbf{0}\ \text{in } B,\qquad
  \bsig\cdot\rhat = \vt\ \text{on } r=a,\qquad
  \bsig = \lambda(\nabla\cdot\vu)\,\mathbf{I}
        + \mu\bigl(\nabla\vu+\nabla\vu^{T}\bigr),
  \label{eq:e1-problem}
\end{equation}
with displacement and traction continuous across every interface (welded).
In the separation \eqref{eq:separation}, the traction on a sphere is
$\bsig\cdot\rhat = T_U\,Y\,\rhat + T_V\,\gradone Y$, with
\begin{equation}
  T_U = \lambda\Bigl(U' + \frac{2U - l(l+1)V}{r}\Bigr) + 2\mu U',
  \qquad
  T_V = \mu\Bigl(V' + \frac{U-V}{r}\Bigr),
  \label{eq:tractions}
\end{equation}
so that at every degree the problem is a two-point boundary-value problem
for $(U,V)$, regular at the centre, with $(U,V,T_U,T_V)$ continuous at the
interfaces and $(T_U,T_V)$ given at the surface. At degree 0, $\gradone Y$
vanishes and only $U$ and $T_U$ remain. This is the bottom rung of the
hierarchy: it tests the elastic operator, the layering and the surface
analysis with nothing else switched on, and on the piecewise-homogeneous
ladder models it has an exact solution.

\subsection{The traction numbers}
\label{sec:e1-numbers}

The forcings are unit tractions, per unit coefficient of the harmonic:
\begin{align}
  \text{normal (pressing):}&\quad \vt = -Y\,\rhat, &
    T_U(a)&=-1,\ T_V(a)=0; \notag\\
  \text{tangential:}&\quad \vt = -\gradone Y, &
    T_U(a)&=0,\ T_V(a)=-1.
\end{align}
The \emph{traction numbers} are the surface values of the response:
\begin{align}
  h_u &= U(a),\quad l_u = V(a) && \text{under the normal traction}; \notag\\
  h_v &= U(a),\quad l_v = V(a) && \text{under the tangential traction}.
\end{align}
They are displacements per unit traction, metres per pascal, and by
\S\ref{sec:harmonics} independent of $m$ and of the harmonics'
normalisation. At degree 0 the tangential channel does not exist and
$l_u=h_v=l_v=0$.

The names are those of the generalised Love numbers, whose ``force''
(subscript $u$) and tangential ($v$) channels are the responses to a unit
surface \emph{density}. A surface density $\sigma$ on the gravitating problem
presses on the surface with the traction $g\sigma$, and \code{pyslfp}'s load
force is built as exactly $g$ times its unit traction (likewise the
tangential forcing), so on the gravitating problem the traction numbers are
the corresponding generalised Love numbers divided by the surface gravity
$g$. Both spectral codes test this tie to $10^{-13}$. The traction numbers
are therefore the channels of the generalised Love numbers that survive
when gravity is switched off, which is what lets one forcing serve E1 and
E2: \code{pyslfp}'s \code{traction\_numbers(model, lmax,
gravitating=True)} solves the full problem under the same unit tractions.

\subsection{Solvability, degree one and degree zero}
\label{sec:e1-solvability}

In weak form, \eqref{eq:e1-problem} asks for $\vu$ with
\begin{equation}
  \int_B \bsig(\vu):\nabla\vw\dV = \int_{\partial B}\vt\cdot\vw\dS
  \quad\text{for every } \vw .
\end{equation}
The left side vanishes for every rigid motion $\vw$ (a translation or an
infinitesimal rotation), so a solution exists only if the traction does no
work on the rigid motions, and it is then unique up to one. A translation
$\mathbf{c}$ is spheroidal of degree 1:
$\mathbf{c} = (\mathbf{c}\cdot\rhat)\,\rhat + \gradone(\mathbf{c}\cdot\rhat)$
with $\mathbf{c}\cdot\rhat$ a degree-1 harmonic. A rotation is toroidal.
Hence, by orthogonality of the harmonics:
\begin{itemize}[leftmargin=*]
\item at $l\ne1$, every spheroidal traction is orthogonal to every rigid
  motion; the problem is solvable, and since rigid motions have no
  spheroidal component of degree $l\ne1$, $U$ and $V$ are unique;
\item at $l=1$, the normal traction $-Y_{1m}\rhat$ does the work
  $-a^2\int Y_{1m}(\mathbf{c}\cdot\rhat)\,\mathrm{d}\Omega\ne0$ on the
  translation along its axis, and the tangential one $-\gradone Y_{1m}$
  twice that: both carry a net force, and there is no static solution.
\end{itemize}
Degree 1 is therefore excluded from E1 and written \code{nan} by every code;
it returns, with gravity, in E2. The exclusion is visible in the closed
form of \S\ref{sec:e1-closed}, which has a pole at $l=1$.

At degree 0 the displacement is radial, and on a homogeneous body the
solution is uniform compression, $\vu = A\vx$, $\bsig = 3\kappa A\,\mathbf
I$ with $\kappa=\lambda+\tfrac23\mu$. A unit pressing traction gives
$3\kappa A=-1$, so
\begin{equation}
  h_u(0) = -\frac{a}{3\kappa},
  \label{eq:degree0}
\end{equation}
which is $-6.4713\times10^{-6}$\,m\,Pa$^{-1}$ on \code{homogeneous}
($\kappa = 3.2817\times10^{11}$\,Pa). Both spectral codes test this closed
form to $10^{-12}$.

\subsection{Reciprocity}
\label{sec:e1-reciprocity}

The elastic operator is self-adjoint, so for two tractions $\vt_1,\vt_2$
with responses $\vu_1,\vu_2$ Betti's reciprocal theorem gives
$\int_{\partial B}\vt_1\cdot\vu_2\dS = \int_{\partial B}\vt_2\cdot\vu_1\dS$.
With $\vt_1=-Y\rhat$ and $\vt_2=-\gradone Y$ at one degree,
\begin{equation}
  -a^2 h_v\!\int Y^2\,\mathrm{d}\Omega
  = -a^2 l_u\!\int |\gradone Y|^2\,\mathrm{d}\Omega
  = -a^2\, l(l+1)\, l_u\!\int Y^2\,\mathrm{d}\Omega ,
  \qquad\text{so}\qquad
  h_v = l(l+1)\,l_u ,
  \label{eq:reciprocity}
\end{equation}
exactly, whatever the layering. A symmetric Galerkin discretisation whose
loads and surface read-outs are the same functionals keeps it exactly, so in
the spectral codes it holds to solver precision, and the residual
$|h_v - l(l+1)l_u|/\max(|h_v|,|l(l+1)l_u|)$ is reported by every leg as a
free internal check. It also halves what has to be compared: $h_v$ carries
no information that $l_u$ does not.

\subsection{The exact reference}
\label{sec:e1-exact}

\path{benchmarks/common/analytic.py} solves E1 exactly on models whose
layers are solid and homogeneous (it samples $\rho$, $v_P$, $v_S$ in each
layer and refuses a layer that is not constant to rounding). The
construction follows.

\paragraph{Harmonic potentials.}
In a homogeneous layer, Navier's equation
$(\lambda+\mu)\nabla(\nabla\cdot\vu)+\mu\nabla^2\vu=\mathbf 0$ has, at
each degree, solutions built from the two harmonic functions
$\chi = r^{n}Y$ with $n=l$ (regular at the centre) or $n=-(l+1)$
(singular); both have $n(n+1)=l(l+1)$, and $\chi$ is homogeneous of degree
$n$, so $(\vx\cdot\nabla)\chi = n\chi$.

\paragraph{Family A: the gradient solution.} $\vu=\nabla\chi$ is
divergence-free and harmonic, so it satisfies Navier's equation. In the
separation,
\begin{equation}
  U = n\,r^{n-1},\qquad V = r^{n-1}.
\end{equation}
Its stress is $2\mu\nabla\nabla\chi$, and since $\nabla\chi$ is homogeneous
of degree $n-1$, $\bsig\cdot\rhat = (2\mu/r)(\vx\cdot\nabla)\nabla\chi =
(2\mu(n-1)/r)\nabla\chi$:
\begin{equation}
  T_U = 2\mu\,n(n-1)\,r^{n-2},\qquad T_V = 2\mu\,(n-1)\,r^{n-2}.
  \label{eq:familyA}
\end{equation}
The same follows from \eqref{eq:tractions}: $T_U = 2\mu U'$ since the
divergence vanishes, and $T_V = \mu\bigl((n-1) - 1 + n\bigr)r^{n-2}$. A
check: at $n=1$ the solution is a uniform translation, and both tractions
vanish.

\paragraph{Family B: the Kelvin combination.} Take
$\vu = c_1\,r^2\nabla\chi + c_2\,\vx\,\chi$. The identities
\begin{equation}
  \nabla\cdot(r^2\nabla\chi) = 2n\chi,\quad
  \nabla\cdot(\vx\chi) = (n+3)\chi,\quad
  \nabla^2(r^2\nabla\chi) = (4n+2)\nabla\chi,\quad
  \nabla^2(\vx\chi) = 2\nabla\chi
\end{equation}
(the third from $\nabla^2(r^2\mathbf f) = 6\mathbf f + 4(\vx\cdot\nabla)
\mathbf f$ for harmonic $\mathbf f$, with $\mathbf f=\nabla\chi$ homogeneous
of degree $n-1$) reduce Navier's operator to
\begin{equation}
  \Bigl\{c_1\bigl[2n(\lambda+\mu)+(4n+2)\mu\bigr]
       + c_2\bigl[(n+3)(\lambda+\mu)+2\mu\bigr]\Bigr\}\nabla\chi ,
\end{equation}
which vanishes identically for
\begin{equation}
  c_1 = (n+3)(\lambda+\mu)+2\mu = (n+3)\lambda+(n+5)\mu,
  \qquad
  c_2 = -\bigl[2n(\lambda+\mu)+(4n+2)\mu\bigr] = -2\bigl[n\lambda+(3n+1)\mu\bigr],
  \label{eq:kelvin}
\end{equation}
the second solution of the classical theory \cite{Love1927}, written
without the division by $c_1$ that the usual normalisation makes. The
division-free form matters: $c_1$ can vanish on the singular side, at
$n = -3-2\mu/(\lambda+\mu)$, for instance at $l=3$ ($n=-4$) in a Poisson
solid ($\lambda=\mu$). The solution is then $c_2\,\vx\chi$, still a solution
and still independent of family A. The two coefficients never vanish
together: that requires $(\lambda+\mu)/\mu\in\{-1,-\tfrac43\}$, excluded
for a stable solid, where $\lambda+\mu=\kappa+\mu/3>0$. In the separation,
\begin{equation}
  U = (c_1 n + c_2)\,r^{n+1},\qquad V = c_1\,r^{n+1},\qquad
  \nabla\cdot\vu = \bigl(2nc_1+(n+3)c_2\bigr)\,r^{n}\,Y,
\end{equation}
\begin{equation}
  T_U = \Bigl[\lambda\bigl(2nc_1+(n+3)c_2\bigr)
              + 2\mu(n+1)(c_1n+c_2)\Bigr]r^{n},
  \qquad
  T_V = \mu\,(2nc_1+c_2)\,r^{n}.
  \label{eq:familyB}
\end{equation}
For the regular solution ($n=l$) these simplify to
$c_1l+c_2=(l+1)[l\lambda+(l-2)\mu]$,
$\nabla\cdot\vu = -2(l+1)(2l+3)\mu\,r^lY$,
$T_U=2\mu(l+1)[(l^2-l-3)\lambda+(l^2-l-2)\mu]\,r^l$ and
$T_V=2\mu[l(l+2)\lambda+(l^2+2l-1)\mu]\,r^l$: family A carries no
dilatation, family B all of it.

The four solutions of degree $l\ge2$ have $U\propto r^{l-1}$, $r^{l+1}$
(regular) and $r^{-l-2}$, $r^{-l}$ (singular), four distinct powers, so they
are independent for every admissible pair of moduli. These powers are the
eigenvalues, $l-1$, $l+1$, $-l-2$, $-l$, of the Euler-type radial system
that the viscoelastic sphere sub-family derives independently
(\code{doc/benchmarks.tex}, ``Sphere benchmarks: problems and exact
references''), which is the basis of the cross-check below. The Navier
equation, the four tractions \eqref{eq:familyA}, \eqref{eq:familyB} and the
degree-0 pair below were checked symbolically in Cartesian coordinates
(sympy) for concrete regular and singular harmonics of degrees 2 and 3
while this note was written.

\paragraph{Degree zero.} With $\gradone Y_{00}=0$ only $U$ survives, and
two of the four solutions collapse. Family A at $n=0$ is the gradient of a
constant, zero; family B at $n=-1$ has $c_1n+c_2=0$, so $U=0$. What remains
is family B at $n=0$, $U\propto r$, and family A at $n=-1$, $U\propto
r^{-2}$:
\begin{equation}
  U = r:\ \ T_U = 3\lambda+2\mu = 3\kappa;
  \qquad
  U = r^{-2}:\ \ T_U = -4\mu\,r^{-3}.
\end{equation}
The first is uniform compression, and with $T_U(a)=-1$ it gives
\eqref{eq:degree0}; the second is the divergence-free field of a
centre of dilatation.

\paragraph{The welded layer system.} Let the layers be $[r_{j-1},r_j]$,
$j=1,\dots,N$, with $r_0=0$ and $r_N=a$, and write the state
$\mathbf s=(U,V,T_U,T_V)^T$ in layer $j$ as $\mathbf s=\mathsf B_j(r)\,
\mathbf c_j$, the columns of $\mathsf B_j$ being the four solutions with
the layer's moduli. Regularity keeps only the two regular columns in the
innermost layer. The interface and surface conditions are
\begin{equation}
  \mathsf B_j(r_j)\,\mathbf c_j = \mathsf B_{j+1}(r_j)\,\mathbf c_{j+1}
  \quad (j=1,\dots,N-1),
  \qquad
  \bigl[\mathsf B_N(a)\,\mathbf c_N\bigr]_{T_U,T_V} = (-1,0)\ \text{or}\ (0,-1),
\end{equation}
$4(N-1)+2$ equations for $2+4(N-1)$ unknowns (at degree 0, with the state
$(U,T_U)$, $2(N-1)+1$ for $1+2(N-1)$). The traction numbers are the
$(U,V)$ rows of $\mathsf B_N(a)\,\mathbf c_N$. The state is propagated
across layer $j$ by the matrix $\mathsf B_j(r_j)\,\mathsf B_j(r_{j-1})^{-1}$,
which in the Euler form is $(r_j/r_{j-1})^{M}$; \code{analytic.py} does not
form these products but assembles and solves the block-bidiagonal system in
all the coefficients at once, so growing and decaying powers are never
multiplied together.

\paragraph{Layer-scaled radii.} In layer $j$ the basis is evaluated at
$x=r/r_j$, the radius scaled by the layer's top. Each column changes by a
constant factor $r_j^{p}$, absorbed into $\mathbf c_j$, and the traction
rows, which carry one radial derivative, are divided by $r_j$ to stay
physical. This is necessary, not cosmetic: in metres, $r^{l}$ overflows
double precision near $l\approx45$ at $r\approx6.4\times10^6$\,m. Scaled,
the regular columns are of order one at the top of their layer and smaller
below it, and the singular columns grow towards the layer's bottom like
$(r_j/r_{j-1})^{l+2}$, a factor of about $4\times10^3$ at $l=32$ across the
upper layer of \code{two\_solid}. The reference is used to $l=32$ here; its
reciprocity residual stays at $\sim10^{-14}$ there
(\S\ref{sec:e1-results}).

\paragraph{The homogeneous sphere in closed form.}
\label{sec:e1-closed}
On one homogeneous layer the system is $2\times2$, and solving it
(symbolically) gives, for $l\ge2$, with
$\Delta_l = (2l^2+4l+3)\lambda + 2(l^2+l+1)\mu$,
\begin{align}
  h_u &= -\frac{a\,\bigl[\,l(2l+1)\lambda + 2(2l^2-1)\mu\,\bigr]}{2\mu(l-1)\,\Delta_l},
  &
  l_u &= -\frac{a\,\bigl[\,3\lambda - 2(l-2)\mu\,\bigr]}{2\mu(l-1)\,\Delta_l},
  \notag\\
  h_v &= l(l+1)\,l_u,
  &
  l_v &= -\frac{a\,\bigl[\,(2l^2+l+3)\lambda + 2(2l^2-l+1)\mu\,\bigr]}{2\mu(l-1)\,\Delta_l}.
  \label{eq:closed}
\end{align}
These agree with \code{analytic.py} on \code{homogeneous} to
$1.6\times10^{-14}$ of the degree scale for $l=2,\dots,32$. They carry the
structure described above and three independent checks:
\begin{itemize}[leftmargin=*]
\item \emph{Reciprocity} \eqref{eq:reciprocity} holds identically, and
  the factor $l-1$ in the denominator is the degree-1 pole of
  \S\ref{sec:e1-solvability}.
\item \emph{The half-space limit.} At high degree, $h_u$ and $l_v$ both
  tend to $-a(\lambda+2\mu)/\bigl(2\mu(\lambda+\mu)\,l\bigr) =
  -(1-\nu)a/(\mu l)$, the compliance $(1-\nu)/(\mu k)$ of an elastic
  half-space at wavenumber $k=l/a$ ($\nu$ Poisson's ratio); on
  \code{homogeneous} at $l=32$ the limit differs from $h_u$ by about
  1\,\%, as an $O(1/l)$ correction should.
\item \emph{Kelvin's rigid sphere.} In the incompressible limit
  $\lambda\to\infty$, $h_u(2)\to -5a/(19\mu)$. In an incompressible
  homogeneous body a harmonic tidal potential $W=(r/a)^2W_aY$ acts only
  through the pressure: absorbing $\rho W$ into it turns the body force into
  a surface traction $\rho W_a Y\rhat$ that pulls. The surface uplift is then
  $5\rho W_a a/(19\mu)$, and $h_2 = gU(a)/W_a = 5\rho g a/(19\mu)$, the
  large-rigidity limit of Kelvin's $h_2 = \tfrac52/(1+19\mu/(2\rho g a))$
  \cite{Love1927}.
\end{itemize}
The formula also shows why the degree-scale measure of
\S\ref{sec:measure} is needed: $l_u$ changes sign at
$l = 2+3\lambda/(2\mu)$, which is $3.96$ on \code{homogeneous}
($\lambda/\mu=1.306$).

\paragraph{A second exact reference.} The viscoelastic sphere sub-family
computes the same per-degree elastic response by a different route: the
Euler system $r\,\mathrm d\mathbf Y/\mathrm dr = M\mathbf Y$ for
$\mathbf Y=(U,V,rT_U,rT_V)$, its regular eigenvectors in the central layer,
and the matrix-exponential propagator $(r_j/r_{j-1})^M$ across the others
(\path{benchmarks/viscoelastic/sphere/reference.py}, derived with sympy by
\path{derive_euler.py}). Its Euler matrix, with real moduli and propagated
as its transfer function does, gives on the E1 ladder models a $U(a),V(a)$
per unit pressure that agrees with \code{analytic.py}'s
$h_u,l_u$ to $3.8\times10^{-15}$ (\code{homogeneous}), $1.8\times10^{-12}$
(\code{homogeneous\_lithosphere}) and $2.5\times10^{-12}$
(\code{two\_solid}), relative to $\max(|h_u|,|l_u|)$ per degree, over
$l=2,\dots,32$. The two references share no code and are derived
differently (explicit potentials and a global solve; the Euler matrix and
propagator products), so E1's reference is itself checked. The check was
made by an ad hoc script while this note was written and is not yet one of
the family's scripts.

\subsection{The legs: what is implemented where}
\label{sec:e1-legs}

\begin{center}
\small
\begin{tabular}{@{}lP{6.3cm}P{3.8cm}l@{}}
\toprule
leg & what computes the traction numbers & gated by & status \\
\midrule
exact & \path{AdGIA/benchmarks/common/analytic.py}; written as exchange
files by \path{elastic/nongrav/reference.py} (\code{e1\_<model>\_analytic.txt}) &
the checks of \S\ref{sec:e1-exact} & done \\
\code{pyslfp} & \code{DegreeSystem(..., gravitating=False)} and
\code{traction\_numbers(model, lmax, gravitating=False)}
(\path{pyslfp/love_numbers/assembly.py}, \path{solver.py}); exchange files by
\path{AdGIA/benchmarks/elastic/nongrav/pyslfp_reference.py} &
nine new tests, the 262-test suite green & done \\
\code{sl3dpp} & \code{DegreeOperator} with
\code{Options::gravitating = false}, and
\code{ComputeTractionNumbers(material, lmax, gravitating)}; exchange files
by the \code{traction\_numbers} driver (\path{sl3dpp/benchmarks/}) &
\code{TestTractionNumbers} against \code{pyslfp} fixtures; the 118-test
suite green & done \\
\code{AdGIA} & \code{traction\_benchmark}
(\path{benchmarks/elastic/nongrav/traction_benchmark.cpp}): the traction
problem on a ladder case, the surface field analysed into coefficients
(\S\ref{sec:e1-fe}) & refinement against the exact reference; reciprocity
and leakage diagnostics & first results \\
\bottomrule
\end{tabular}
\end{center}

\paragraph{\code{pyslfp}.} The non-gravitating problem uses the released
degree discretisation with the gravity and potential terms dropped: the
unknowns are $U$ and $V$ ($U$ alone at $l=0$), the elastic block of the
bilinear form is unchanged, and the natural forcings are the unit tractions,
$-a^2$ on the surface $U$ test function (normal) and $-a^2l(l+1)$ on the
surface $V$ test function (tangential). The class refuses a fluid layer
(a static non-gravitating fluid has no restoring force) and degree 1. On
the gravitating problem the same forcings give the gravitating traction
numbers, and the existing load forcings are exactly $g$ times them. Among
the new tests: the degree-0 closed form to $10^{-12}$; the exclusion of
degree 1; reciprocity below $10^{-10}$; the tie to the gravitating solver to
$10^{-13}$; and a mesh-independence check that is exact in its own right. On
a homogeneous body the regular solutions are polynomials of degree $l+1$,
which lie in the five-node element space for $l\le3$, so two unrelated
meshes agree to rounding there.

\paragraph{\code{sl3dpp}.} The C++ mirror is exact by design:
\code{DegreeOperator} drops the same terms and degrees of freedom,
\code{Traction()} and \code{TractionTangential()} are the unit tractions
(the full problem's \code{Load(Force)} is $g$ times \code{Traction()}),
and \code{ComputeTractionNumbers} returns the SI surface response with the
same \code{nan} at degree 1. \code{TestTractionNumbers} holds it to
\code{pyslfp}: the fixtures (\path{scripts/make_e1_fixtures.py}, run from the
\code{pyslfp} checkout) are the \code{homogeneous} and \code{two\_solid}
decks of \code{AdGIA}'s ladder at full precision, read back with
planetmodel, and \code{pyslfp}'s traction numbers on the uniform
$l_{\max}=8$ rule, non-gravitating and gravitating; \code{sl3dpp} reads the
same decks on the same rule. The test also checks the degree-0 closed form,
reciprocity, the $g$-tie to \code{ComputeLoveNumbers}, and the refusals.
The \code{traction\_numbers} driver (\code{traction\_numbers <deck> <lmax>
<out.txt> [nongrav|grav]}) writes the exchange file.

\subsection{Results}
\label{sec:e1-results}

All differences are relative to the degree scale \eqref{eq:degreescale}
unless stated, SI throughout.
\begin{center}
\small
\begin{tabular}{@{}P{4.2cm}P{6.1cm}P{4.9cm}@{}}
\toprule
comparison & conditions & worst difference \\
\midrule
\code{sl3dpp} against \code{pyslfp} (identical discrete problem) &
\code{homogeneous}, \code{two\_solid} decks; uniform $l_{\max}=8$ rule,
five GLL nodes; non-gravitating and gravitating; degrees 0--8 (1 excluded
without gravity) & within the gate, $10^{-10}$ \\ \addlinespace
\code{pyslfp} against the exact reference & graded mesh, $l_{\max}=16$,
five GLL nodes & $9.4\times10^{-11}$ (\code{homogeneous}),
$1.8\times10^{-11}$ (\code{homogeneous\_lithosphere}), $1.7\times10^{-11}$
(\code{two\_solid}) \\ \addlinespace
\code{pyslfp} fixtures against the exact reference & the committed
\code{sl3dpp} fixtures: uniform $l_{\max}=8$ rule; measured for this note &
$1.2\times10^{-10}$ (\code{homogeneous}), $6.5\times10^{-11}$
(\code{two\_solid}) \\ \addlinespace
\code{AdGIA} against the exact reference & \code{two\_solid}, degrees
0--8; $h=0.3$, order 2, and $h=0.15$, order 3 (\S\ref{sec:e1-fe}) & $h_u$:
$6.6\times10^{-2}$, $1.6\times10^{-3}$; all channels:
$6.3\times10^{-1}$, $1.4\times10^{-2}$ (in $h_v$, at $l=8$) \\ \addlinespace
reciprocity, \code{AdGIA} & the same two runs & $1.2\times10^{-2}$,
$8.5\times10^{-8}$ \\ \addlinespace
reciprocity, exact reference & $l=2,\dots,32$, the three models &
$\le4\times10^{-14}$ (relative to $\max(|h_v|,|l(l+1)l_u|)$) \\ \addlinespace
reciprocity, spectral codes & as tested & $<10^{-10}$ \\ \addlinespace
degree 0, \eqref{eq:degree0} & homogeneous body; both spectral codes &
$10^{-12}$ (tested); exact reference to rounding \\ \addlinespace
the two exact references & Euler propagator against \code{analytic.py},
$l=2,\dots,32$, normal channel & $3.8\times10^{-15}$, $1.8\times10^{-12}$,
$2.5\times10^{-12}$ \\ \addlinespace
closed form \eqref{eq:closed} & against \code{analytic.py},
\code{homogeneous}, $l=2,\dots,32$ & $1.6\times10^{-14}$ \\ \addlinespace
\code{homogeneous\_lithosphere} against \code{homogeneous} & exact
reference, $l\le32$: one body, cut by a welded interface with no contrast &
$2.6\times10^{-15}$ \\
\bottomrule
\end{tabular}
\end{center}

The first row is tolerance tier (i), the next two tier (ii) and the
\code{AdGIA} rows tier (iii) of \S\ref{sec:tolerances}. Together they place \code{sl3dpp} within about
$2\times10^{-10}$ of the exact numbers on \code{homogeneous} and
\code{two\_solid} (by the triangle inequality, through the fixtures; an
\code{sl3dpp} exchange file has not yet been compared with the exact
reference directly). \code{pyslfp}'s differences from the exact reference
depend on the mesh rule, as tier (ii) expects: on \code{homogeneous} the
graded mesh at $l_{\max}=16$ and the uniform rule at $l_{\max}=8$ give
$9.4\times10^{-11}$ and $1.2\times10^{-10}$. The degree-0 agreement of the
gravitating fixtures also verifies, for solid models, that \code{sl3dpp}
treats degree 0 as \code{pyslfp} does (\S\ref{sec:degrees01}). The
finite-element numbers are discussed in \S\ref{sec:e1-fe}.

\paragraph{Reference values.} For orientation, the exact traction numbers,
in $10^{-6}$\,m\,Pa$^{-1}$ (\code{homogeneous\_lithosphere} is identical to
\code{homogeneous}):
\begin{center}
\small
\begin{tabular}{@{}r rrrr c rrrr@{}}
\toprule
& \multicolumn{4}{c}{\code{homogeneous}} && \multicolumn{4}{c}{\code{two\_solid}} \\
\cmidrule{2-5}\cmidrule{7-10}
$l$ & $h_u$ & $l_u$ & $h_v$ & $l_v$ && $h_u$ & $l_u$ & $h_v$ & $l_v$ \\
\midrule
0 & $-6.4713$ & 0 & 0 & 0 && $-7.9708$ & 0 & 0 & 0 \\
2 & $-13.349$ & $-1.9326$ & $-11.596$ & $-15.281$ && $-12.965$ & $-1.9089$ & $-11.453$ & $-27.502$ \\
3 & $-8.5106$ & $-0.26567$ & $-3.1880$ & $-8.7762$ && $-9.8339$ & $-0.56925$ & $-6.8310$ & $-18.614$ \\
4 & $-6.4065$ & $0.0048571$ & $0.097141$ & $-6.4016$ && $-8.5523$ & $-0.23353$ & $-4.6705$ & $-14.055$ \\
8 & $-3.2897$ & $0.061608$ & $4.4358$ & $-3.2281$ && $-6.2731$ & $0.052111$ & $3.7520$ & $-7.0881$ \\
16 & $-1.6778$ & $0.023609$ & $6.4216$ & $-1.6542$ && $-3.7719$ & $0.046307$ & $12.595$ & $-3.7451$ \\
32 & $-0.84812$ & $0.0069845$ & $7.3756$ & $-0.84113$ && $-1.9282$ & $0.014496$ & $15.307$ & $-1.9137$ \\
\bottomrule
\end{tabular}
\end{center}

\paragraph{What gravity changes.} The gravitating fixtures show the size of
the step to E2 under the same forcing (\code{pyslfp}, uniform
$l_{\max}=8$ rule). The ratio of the gravitating to the non-gravitating
$h_u$ is $1.43$ at degree 0 and $0.935$--$0.976$ over degrees 2--8 on
\code{homogeneous} (for $l_v$, $0.997$--$1.049$), and $1.20$ at degree 0
and $0.974$--$0.984$ over degrees 2--8 on \code{two\_solid}: gravity
softens the degree-0 response and stiffens the radial response of the
others by a few per cent. Each code's non-gravitating leg must be right
to far better than this before E2 can say anything about gravity.

\subsection{Reproducing}
\label{sec:e1-reproduce}

From the build tree, \path{<build>/benchmarks/elastic/nongrav/}
(\path{benchmarks/elastic/nongrav/README.md} has the details):
\begin{verbatim}
./reference --lmax 32 --out runs                    # the exact reference
(cd ~/dev/pyslfp && poetry run python \
    <repo>/benchmarks/elastic/nongrav/pyslfp_reference.py --lmax 16 \
    --out <build>/benchmarks/elastic/nongrav/runs)  # pyslfp, graded mesh
./compare runs/e1_two_solid_analytic.txt runs/e1_two_solid_pyslfp.txt \
    --tol 1e-8 --verbose
\end{verbatim}
\code{pyslfp\_reference.py --gravitating} writes the gravitating twin under
the same forcing (\path{e1_<model>_pyslfp_grav.txt}). For \code{sl3dpp}, the
decks come from \code{python <repo>/benchmarks/common/decks.py homogeneous
two\_solid --out <dir>} and the exchange file from \code{traction\_numbers
<dir>/<model>.elastic.deck <lmax> <out.txt>}; its gate is
\code{ctest -R TractionNumbers} in the \code{sl3dpp} build. For \code{AdGIA},
from the same directory, a case and a run (the finer of the two of
\S\ref{sec:e1-fe}):
\begin{verbatim}
../../love_numbers/make_case two_solid --h 0.15 --out cases/two_solid_h015
mpiexec -np 8 ../../bin/traction_benchmark \
    -c cases/two_solid_h015/case.json -lmax 8 -o 3 \
    -out e1_two_solid_adgia_h015_o3.txt
./compare runs/e1_two_solid_analytic.txt e1_two_solid_adgia_h015_o3.txt
\end{verbatim}

\subsection{The finite-element leg}
\label{sec:e1-fe}

\code{AdGIA}'s leg is the one E1 exists for. \code{traction\_benchmark}
(\path{benchmarks/elastic/nongrav/traction_benchmark.cpp}, built with the
other drivers) takes a case written by \code{make\_case.py} for a ladder
model: its mesh, $\kappa$ and $\mu$, nothing about the model set in the
driver, and a refusal if the case has a fluid region. It solves
\eqref{eq:e1-problem} with \code{LinearQuasiStaticTractionProblem} on the
body alone (the buffer shell of the gravitating cases is left out), with no
gravity and no potential and the rigid modes projected out in the solver.
For each degree, order 0, the surface traction is $-Y_{l0}\rhat$ and, for
$l\ge2$, $-\gradone Y_{l0}$, applied one at a time; the displacement is
analysed into radial and tangential harmonic coefficients on the surface
(\code{BoundaryHarmonicCoefficients}); and the $(l,0)$ coefficients,
converted to SI through the case's scales, are written as an exchange file,
degree 1 as \code{nan}. The driver also reports the reciprocity residual
and \code{spurious}, the largest unforced coefficient relative to the
forced ones. The machinery is not new: the viscoelastic sphere sub-family
solves the same free-body pressure problem on curved layered balls with the
same classes (\S\ref{sec:internal-visco}). What E1 adds is the shared models
in SI, the tangential channel, and a reference that the spectral codes are
held to as well.

\paragraph{First measurements.} Two runs on \code{two\_solid}, degrees
0--8, solver tolerance $10^{-12}$: $h=0.3$ at order 2 (4481 elements,
buffer shell included in the count) and $h=0.15$ at order 3 (14\,697
elements). The driver's commit records them in
\path{benchmarks/elastic/nongrav/README.md} (\code{faddd52}); the
differences below were recomputed for this note with \code{compare.py}
from the two runs' exchange files against the exact reference, with the
degree-scale measure \eqref{eq:degreescale}:
\begin{center}
\small
\begin{tabular}{@{}lcccc@{}}
\toprule
& \multicolumn{2}{c}{$h=0.3$, order 2} & \multicolumn{2}{c}{$h=0.15$, order 3} \\
\cmidrule(lr){2-3}\cmidrule(l){4-5}
channel & $l=2$ & worst ($l=8$) & $l=2$ & worst ($l=8$) \\
\midrule
$h_u$ & $5.6\times10^{-4}$ & $6.6\times10^{-2}$ & $2.1\times10^{-6}$ & $1.6\times10^{-3}$ \\
$l_u$ & $7.0\times10^{-5}$ & $8.7\times10^{-3}$ & $2.0\times10^{-6}$ & $1.9\times10^{-4}$ \\
$h_v$ & $4.2\times10^{-4}$ & $6.3\times10^{-1}$ & $1.2\times10^{-5}$ & $1.4\times10^{-2}$ \\
$l_v$ & $1.2\times10^{-3}$ & $1.2\times10^{-1}$ & $7.0\times10^{-6}$ & $2.1\times10^{-3}$ \\
\midrule
$h_u$, degree 0 & \multicolumn{2}{c}{$7.3\times10^{-5}$} & \multicolumn{2}{c}{$4.3\times10^{-6}$} \\
reciprocity & $2.2\times10^{-6}$ & $1.2\times10^{-2}$ & $1.4\times10^{-11}$ & $8.5\times10^{-8}$ ($l=6$) \\
\bottomrule
\end{tabular}
\end{center}
Every channel's error grows with degree, steeply on the coarse mesh, where
degree 8 is at the resolution limit (as in the Love-number family at the
same $h$), and the finer, higher-order run lowers it by factors of 35 to
270 across the entries of the table. The largest entries are in
$h_v$, and they are not new information: by reciprocity $h_v\approx
l(l+1)\,l_u$, so its error is the absolute error of the small channel $l_u$
multiplied by $l(l+1)$, 72 at degree 8 ($8.7\times10^{-3}\times72\approx
6.3\times10^{-1}$; $1.9\times10^{-4}\times72\approx1.4\times10^{-2}$).
Reciprocity, exact in the continuum and to solver precision in the spectral
codes, holds in the finite-element leg at the discretisation level, to
$1.4\times10^{-11}$ at degree 2 on the finer run. The two runs change $h$
and the order together: they show tolerance tier (iii) behaving as
expected, but they are not a convergence study.

\paragraph{Next.} The refinement study: $h$- and order-ladders on
\code{homogeneous}, \code{homogeneous\_lithosphere} and \code{two\_solid},
with rates. On \code{homogeneous\_lithosphere} the numbers must also match
\code{homogeneous}'s to the discretisation level although the mesh differs,
since the two are the same body.

Open points of E1 beyond the finite-element leg: an \code{sl3dpp} exchange
file compared with the exact reference directly; the cross-check of the two
exact references made a script of the family; and E1 as a campaign stage.

%-----------------------------------------------------------------------------
\section{Tier E2: gravitating elastic, solid models}
\label{sec:e2}
%-----------------------------------------------------------------------------

\status{21 October 2026}{Begun. The spectral pair's gravitating traction
numbers are gated against each other (the E1 gate). \code{AdGIA}'s
Love-number family against \code{pyslfp} is in place, and the first bridges
that bring it under the programme's conventions --- exchange files of the
conventional Love numbers from \code{love\_benchmark}'s results and from
\code{pyslfp} (\path{benchmarks/elastic/grav/}) --- are being written. The
degree-1 frame work is pending.}

\paragraph{The problem.} The full linearised, quasi-static, self-gravitating
elastic response of solid spherical models to surface loads and tidal
potentials, with the perturbation of the potential solved inside and
outside the body. The model ladder: \code{homogeneous}, \code{two\_solid},
\code{linear\_solid}. The gravity split is two-way, no gravity (E1) or full
gravity (E2); an intermediate variant (background gravity with the
potential perturbation dropped) is added only if an issue arises that it
would help isolate.

\paragraph{The compared quantities.} The conventional load Love numbers
$h'_l,l'_l,k'_l$ and tidal numbers $h_l,l_l,k_l$ (dimensionless; the tidal
ones from degree 2), which every code computes; the gravitating traction
numbers (the same forcing as E1, so that the pair isolates gravity); and
the generalised Love numbers channel by channel where a code exposes them.
Degree 1 is compared in the centre-of-mass frame (\S\ref{sec:degrees01});
degree 0 is included on solid models.

\paragraph{References and legs.} \code{pyslfp} is the reference; no closed
form exists for a compressible, gravitating, layered body. The spectral
pair already meets tolerance tier (i) on the gravitating traction numbers
(\code{homogeneous} and \code{two\_solid}, $l_{\max}=8$, degrees 0--8
including 1). \code{AdGIA}'s leg is the Love-number family
(\S\ref{sec:internal-elastic}), which solves exactly these problems against
\code{pyslfp} on its own case meshes; it is to be revisited, brought into the
exchange format and SI, and its degree-1 frame settled against
\code{pyslfp}'s. The two \code{AdGIA} solver routes benchmarked are the
Dahlen (mixed Eulerian) path and the Maxwell-core route; the other
formulations are kept for what they show but do not enter the benchmark
matrix.

\paragraph{Open points.} The degree-1 frame; the choice between the graded
and uniform spectral mesh rules for the reference files; a gravitating
analogue of the E1 FE driver or an extension of it.

%-----------------------------------------------------------------------------
\section{Tier E3: gravitating elastic, fluid--solid models}
\label{sec:e3}
%-----------------------------------------------------------------------------

\status{21 October 2026}{Pending as a tier. \code{AdGIA}'s Love-number
family covers these models against \code{pyslfp}.}

\paragraph{The problem.} E2 with fluid regions: a fluid core, then an inner
core, stratified cores and PREM-grade structure (\code{fluid\_core},
\code{aw\_core}, \code{inner\_core}, \code{stratified\_core},
\code{earth\_like}, \code{prem\_4}). There is no non-gravitating fluid--solid
tier: a static fluid without gravity has no restoring force.

\paragraph{Degree zero.} The place where the codes' treatments of degree 0
differ (\S\ref{sec:degrees01}): \code{pyslfp} solves the compressible fluid;
\code{sl3dpp} is to be verified to do the same with a fluid layer;
\code{AdGIA}'s Dahlen path differs by design at $l=0$ and is shown, not
benchmarked, there (a relative $h'_0$ error of about 10 on
\code{fluid\_core}, \code{doc/benchmarks.tex}, ``Degree zero with a fluid
layer''), while the Maxwell-core route is benchmarked at every degree.

\paragraph{What is known.} The Love-number family's conclusions on these
models (\S\ref{sec:internal-elastic}), including the caveat that a
non-neutrally stratified core separates the static formulations from the
secular response at $O(N^2)$, so that \code{aw\_core}, neutral by
construction, is the clean comparison model.

%-----------------------------------------------------------------------------
\section{Beside the elastic tiers: \code{AdGIA}'s internal families}
\label{sec:internal-elastic}
%-----------------------------------------------------------------------------

These families test \code{AdGIA} alone, against \code{pyslfp} as an
independent radial reference or against exact identities. Their problems,
settings and full measurements are in \code{doc/benchmarks.tex}; what they
have established, in brief:

\paragraph{Love numbers (\code{love\_numbers/}).} Load and tidal Love
numbers of layered, self-gravitating elastic bodies, every formulation
against \code{pyslfp}, with the response also compared as fields under a cap
load (\code{doc/benchmarks.tex}, ``The Love-number family''). Every
formulation agrees with \code{pyslfp} at the discretisation level on its
domain of validity: at $h=0.2$ the median relative error over degrees 1--5
is $5\times10^{-3}$--$1.4\times10^{-2}$ at order 2 and
$6\times10^{-5}$--$8\times10^{-4}$ at order 3, for every model from
\code{homogeneous} to \code{prem\_4}; the spread between formulations at
coarse resolution is mesh error, not physics. For loading at $l\ge1$ on
spherical models the Dahlen path is the cheapest and as accurate as any.
The dearer formulations buy what it cannot represent --- the compressible
fluid at degree 0, the slipping interface, the mapped machinery --- but not
non-neutral core stratification, where the general method is the Maxwell
relaxation of the fluid. The standard unmeshed-core CMB condition
(\code{uniform}) is safe for loading at $l\ge2$ ($h'$ within $0.1\,\%$,
$l'_2$ and $k'_2$ about $0.3\,\%$) and about ten times worse on tides;
buoyancy alone (\code{winkler}) is safe for neither. The fields show what the
surface numbers do not: on models with a solid inner core the order-2
displacement error is large ($0.49$ on \code{prem\_4} at $h=0.2$) where the
Love numbers are good to a per cent, through the weakly restored
translation of the inner core, and order 3 removes it. Open: the gauged
method's degree-1 interior potential, and the degree-1 frame against
\code{pyslfp}.

\paragraph{Relabelling (\code{relabelling/}).} The mapped (aspherical)
machinery, tested with the spherical physics unchanged by describing the
same problem in relabelled coordinates \cite{AlAttarCrawford2016}
(\code{doc/benchmarks.tex}, ``The relabelling family''). The discrete
change-of-variables identity holds for the welded classes to
$\sim10^{-7}$ ($A=0.02$, $h=0.3$, order 2), and on the sea-level legs to
$5\times10^{-7}$ and below; the mapped Love numbers agree with the unmapped
ones at the level of the errors themselves; an aspherical reference body
gives field errors independent of the shape amplitude. Through the
slipping interface the identity does not hold: the mapped slipping forms are
not certified.

\paragraph{Perturbation (\code{perturbation/}).} Derivatives of the response
with respect to an interface radius, the mapped assembly under a
discontinuous deformation gradient against central differences of
\code{pyslfp} solves (\code{doc/benchmarks.tex}, ``The perturbation
family''). On \code{fluid\_core} the degree-0 derivative agrees to
$0.8\,\%$ at order 2 and $3\times10^{-3}\,\%$ at order 3; at degrees 2 and 3
the discrepancies of 1--16\,\% are of the order of the errors of the
numbers themselves at $h=0.3$, and on \code{two\_solid} 3--17\,\%: a
consistency check at this resolution rather than a sharp verification.

\paragraph{The disc (\code{disc/}).} The two-dimensional laboratory of the
fluid treatments (\path{benchmarks/disc/README.md} and
\code{doc/planning/no\_gravity\_analytics.md}; it has no section in
\code{doc/benchmarks.tex}). Without gravity a fluid disc in a solid annulus
has closed-form solutions and every tangential slip is removable, so
welded and slipping solves must coincide exactly, which tests the constraint
machinery alone. With gravity, a per-degree radial reference with welded,
free-slip and Dahlen switches measures what can be separated: on the
Adams--Williamson disc core the three agree to about $10^{-10}$; on a
uniform, non-neutral disc core the welded and free-slip solves drift within
an envelope of about $0.3\,\%$ under refinement
(\code{doc/benchmarks.tex}, ``The formulations compared'').

%=============================================================================
\part{Viscoelastic loading}
\label{part:visco}
%=============================================================================

\section{Tier V: viscoelastic loading without the sea-level equation}
\label{sec:v}

\status{21 October 2026}{Pending. The ingredients exist in each code
separately: \code{AdGIA}'s viscoelastic family
(\S\ref{sec:internal-visco}); \code{pyslfp}'s time-domain viscoelastic
module; \code{sl3dpp}'s viscoelastic stepping gated against \code{pyslfp}
fixtures.}

\paragraph{The problem.} The response of spherical models to prescribed
surface-load histories, with no sea-level equation, in all three codes,
with the same splits as the elastic tiers: the model ladder; gravity off and
on; rotational feedback off and on. The three codes share the
exponential-trapezoid step weights, so with matched steps a comparison
measures the spatial discretisation.

\paragraph{References.} The spectral codes, against each other at tier
(i) where they solve the identical discrete problem. On the non-gravitating
rung the E1 reference extends exactly: by the correspondence principle the
transfer functions $U(a;s)$, $V(a;s)$ are rational in the Laplace variable,
which is how the viscoelastic sphere sub-family's exact histories are built.
Semi-analytical references (modal sums, AAA with Talbot, Gaver--Stehfest)
are kept as secondary.

\paragraph{The lateral pair (V-lat).} \code{AdGIA} and \code{sl3dpp} are the
two codes with laterally varying viscosity. With a radial elastic structure
and viscosity $\bar\eta(r)\,e^{\delta(r,\theta,\varphi)}$, the two can be
compared for loading without sea level, the viscosity field exchanged as
netCDF. This needs \code{sl3dpp}'s functionality exposed through public
drivers.

\section{Beside tier V: the viscoelastic family}
\label{sec:internal-visco}

\code{AdGIA}'s viscoelastic family (\path{viscoelastic/},
\code{doc/benchmarks.tex}, ``The viscoelastic family'') has four
sub-families. The box and sphere sub-families are non-gravitating and
solid-only, with references exact to about $10^{-10}$, so any mismatch is
the code's: on them every integrator reaches its design order on every body
it is stable on, with one finding beyond the design orders --- the
exponential trapezoid loses half an order (1.5 against 2) under stress
control on a stiff multi-branch body; and order-2 elements converge at
about fourth order in the surface coefficients on curved spherical meshes.
The sphere sub-family's free layered balls under a degree-$l$ pressure are
the viscoelastic version of E1, and its laterally varying viscosity runs
test the integrators against a much smaller step. The stepping survey
compares the schemes on one beam. The Maxwell Love-number histories of
gravitating models (\path{viscoelastic/love/}) are compared with
\code{pyslfp} through the correspondence principle and Gaver--Stehfest
inversion: history errors of $0.5$--$1.1\times10^{-2}$ in $h'$ at order 2 on
\code{fluid\_core} ($h=0.3$) and $1.7$--$3.2\times10^{-4}$ at order 3, the
errors spatial, not temporal. The choice of integrator these support is
described in \code{doc/viscoelasticity.md}.

%=============================================================================
\part{Sea level}
\label{part:sea}
%=============================================================================

\section{Tier S: sea level}
\label{sec:s}

\status{21 October 2026}{Pending, and to start only when E and V agree. An
\code{AdGIA}--\code{pyslfp} comparison (\S\ref{sec:internal-sea}) and
\code{sl3dpp}'s stage-2 sea-level gates against \code{pyslfp} exist
separately.}

\paragraph{The problem.} The sea-level equation on the same model
hierarchy: elastic fingerprints first, three-way, with and without
rotational feedback and shoreline migration; then full viscoelastic GIA
runs; and 3-D viscosity with the two codes that have it. All three codes
solve the generalised sea-level equation of \code{pyslfp}, of which
\code{sl3dpp}'s is a port and \code{AdGIA}'s monolithic treatment a
reformulation.

\paragraph{Open points.} The observables and their observational
yardsticks (\S\ref{sec:open}); the degree-1 frame of the sea-level
comparison, flagged as unreconciled in the existing one.

\section{Beside tier S: the sea-level family}
\label{sec:internal-sea}

\code{AdGIA}'s monolithic sea-level solve (\code{SetWaterLoad}) against
\code{pyslfp}'s pseudo-spectral solution of the same problem, with
rotation, shoreline-migration and timing legs and a real-geography variant
on coastline-refined meshes (\path{benchmarks/sea_level/README.md}; the
family has no section in \code{doc/benchmarks.tex}). At toy resolution
(\code{homogeneous}, $h=0.4$, order 2, $l_{\max}=16$) the ocean-RMS
relative difference is $1.1\,\%$ on the Dahlen path and $1.2\,\%$ on the
welded referential class, the floor of two discretisations at this
resolution rather than a disagreement; the resolved comparison belongs to
the server campaign. Love-number agreement is assumed there and tested by
the Love-number family, which is why the elastic tiers come first.

%=============================================================================
\begin{thebibliography}{99}
%=============================================================================

\bibitem{AlAttarCrawford2016}
D.~Al-Attar \& O.~Crawford, 2016.
Particle relabelling transformations in elastodynamics,
\emph{Geophys.\ J.\ Int.}, 205, 575--593.

\bibitem{Love1927}
A.~E.~H.~Love, 1927.
\emph{A Treatise on the Mathematical Theory of Elasticity}, 4th edn,
Cambridge University Press.

\bibitem{Martinec2018}
Z.~Martinec, V.~Klemann, W.~van der Wal, R.~E.~M.~Riva, G.~Spada, Y.~Sun,
D.~Melini, S.~B.~Kachuck, V.~Barletta, K.~Simon, G.~A \& T.~S.~James, 2018.
A benchmark study of numerical implementations of the sea level equation in
GIA modelling, \emph{Geophys.\ J.\ Int.}, 215, 389--414.

\bibitem{Spada2011}
G.~Spada, V.~R.~Barletta, V.~Klemann, R.~E.~M.~Riva, Z.~Martinec,
P.~Gasperini, B.~Lund, D.~Wolf, L.~L.~A.~Vermeersen \& M.~A.~King, 2011.
A benchmark study for glacial isostatic adjustment codes,
\emph{Geophys.\ J.\ Int.}, 185, 106--132.

\end{thebibliography}

\end{document}
